\documentclass[10pt,a4paper]{amsart}
\usepackage[margin=19mm]{geometry}
\usepackage{amsmath,amssymb,mathtools}
\usepackage[scr=rsfs,cal=euler]{mathalfa}
\usepackage{xfrac}
\usepackage[unicode,hidelinks,bookmarks=false]{hyperref}
\newcommand{\ie}{i.e.\ }
\newcommand{\cf}{cf.\ }
\newcommand{\resp}{respectively\ }
\newcommand{\Subsection}{\subsection}
\providecommand{\nobreakdash}{\nobreak\hbox{-}}
\setlength{\emergencystretch}{2em}
\multlinegap0pt
\usepackage{amssymb}
\usepackage{mathtools}
\newtheorem{lemma}{Lemma}
\usepackage{cleveref}
\crefname{lemma}{Lemma}{Lemmas}
\title{An 18-colour bound for locally irregular decompositions}
\author{Carla Negri Lintzmayer, Guilherme Oliveira Mota, Maycon Sambinelli, Vinicius Fernandes dos Santos}
\date{Source: arXiv:2609.09355v1 (CC BY 4.0)}
\begin{document}
\maketitle

\noindent\emph{Source and licence.} An 18-colour bound for locally irregular decompositions. Version arXiv:2609.09355v1 (CC BY 4.0), distributed by arXiv under the Creative Commons Attribution 4.0 International licence, http://creativecommons.org/licenses/by/4.0/, as recorded in the arXiv licence metadata for this exact version and shown on the abstract page https://arxiv.org/abs/2609.09355v1. Full licence notice: \texttt{licenses/arxiv-2609.09355-CC-BY-4.0.txt}.\par\smallskip

\noindent\textit{Editorial scope.} Counted unit: the article's lemma lem:marked (source0.tex lines 409--415). The article's complete original proof of it is delivered with it (source0.tex lines 416--503, one \texttt{proof} environment of 88 lines). No mathematics is written, repaired or re-derived: the statement and the proof are the article's own lines, in the article's own notation, and no retained line is substituted or re-flowed.  Retained uncounted as the source-local context the proof needs: lines 299--311; lines 341--347; lines 400--403.  Nothing was omitted from any retained span, so the package carries no omission record.  No retained line contains a citation, so the wrapper carries no bibliography and no bibliography entry.  No retained line contains a figure environment, an image-inclusion command or drawing code, so no image is delivered and the package declares no asset dependency.  Editorial formatting only, in addition: an AMS \texttt{amsart} preamble that restates the article's own declarations and macros used by the retained lines, namely the environments \texttt{lemma} on separate counters, together with the standard packages those lines need.  The retained environments print in this document as Lemma 1 in order of appearance, whereas the article prints the same items with the numbering of its own full document.  The complete native source of the article as distributed in the arXiv e-print for this exact version is bundled unchanged as \texttt{sources/209804/source0.tex}.
\begin{lemma}\label{lem:parity-orientation}
  Let~$L$ be a connected loopless multigraph. 
  Then
  \begin{enumerate}
    \renewcommand{\theenumi}{\roman{enumi}}
    \item For every map $\eta \colon V(L) \to \{0,1\}$ and every spanning tree~$T$
    of~$L$, every orientation of $L-E(T)$ extends to an orientation $\vec L$
    of~$L$ satisfying $d^-_{\vec L}(v) \equiv \eta(v) \pmod 2$ for all~$v$ if and
    only if $\sum_{v \in V(L)} \eta(v) \equiv e(L) \pmod 2$.
    \item There is an orientation of $L$ in which all indegrees are even,
    except possibly one, which may have indegree~$1$.
  \end{enumerate}
\end{lemma}

\begin{lemma}\label{lem:lists} 
  Let~$N$ be a bipartite loopless multigraph with finite forbidden sets of
  non-negative integers~$F(v)$ for each vertex~$v$.
  If~$N$ has an orientation~$D$ satisfying $d^-_D(v) \ge |F(v)|$ at every
  vertex~$v$, then~$N$ has an orientation~$O$ satisfying $d^-_O(v) \notin F(v)$
  at every~$v$.
\end{lemma}

\begin{equation}\label{eq:hall}
  \sum_{v \in S} |F(v)| \le e_N(S) + |\partial_N(S)|
    \qquad\text{for every }S \subseteq V(N) \, ,
\end{equation}

\begin{lemma}\label{lem:marked}
  Let~$M$ be a bipartite loopless multigraph, and prescribe a common
  bipartition direction on a matching~$P$ of at most two edges.
  Let~$k$ be even, with $k \ge 2$ if $|P| \le 1$ and $k \ge 4$ otherwise.
  The prescription extends to an orientation whose indegrees are even or~$1$,
  and in which no head of~$P$ has indegree~$k$.
\end{lemma}

\begin{proof}
  Call the heads of~$P$ marked and all other vertices ordinary.
  Put $N := M-P$.
  Once the prescribed arcs of~$P$ are restored, the total indegree of a marked
  head is its indegree in~$N$ plus~$1$, whereas the total indegree of an
  ordinary vertex is its indegree in~$N$.
  Accordingly, for an ordinary vertex~$v$, define
  \[
    F(v) := \{r \in \{3,4,\ldots,d_N(v)\} \colon r \text{ is odd }\} \, ,
  \]
  and, for a marked vertex~$v$, define
  \[
    F(v) := \{r \in \{0,1,\ldots,d_N(v)\} \colon r \text{ is positive and even,
    or } r = k-1\}.
  \]
  Put $\delta(v) := 2|F(v)| - d_N(v)$.
  Since the prescribed edges have a common bipartition direction, all marked
  heads lie in the same bipartition class.

  At an ordinary nonisolated vertex, $\delta(v)$ is~$-1$ when $d_N(v)$ is odd
  and~$-2$ when~$d_N(v)$ is even.
  At a marked head with $d_N(v) \ge k-1$, the value $\delta(v)$ is~$1$ when
  $d_N(v)$ is odd and~$2$ when $d_N(v)$ is even; at any other marked head it
  is~$-1$ or~$0$.
  An isolated ordinary vertex has $\delta(v) = 0$.

  If~\eqref{eq:hall} holds, the preceding Hall criterion and \Cref{lem:lists}
  give an orientation of~$N$ avoiding every~$F(v)$.
  After the prescribed arcs of~$P$ are restored, every total indegree is even
  or~$1$, and no marked head has total indegree~$k$.

  We may therefore assume that~\eqref{eq:hall} fails.
  The difference $\sum_{v \in S} |F(v)| - e_N(S) - |\partial_N(S)|$ is additive
  over the components of~$N[S]$, so there is a set~$S$ such that~$N[S]$ is
  connected and this difference is positive. 
  Since $\sum_{v \in S} d_N(v) = 2e_N(S) + |\partial_N(S)|$, this is equivalent
  to
  \begin{equation}\label{eq:deficiency}
    \sum_{v \in S} \delta(v) - |\partial_N(S)|
      = 2 \left(\sum_{v \in S} |F(v)| - e_N(S) - |\partial_N(S)| \right) \ge 2 \, .
  \end{equation}
  A singleton cannot satisfy~\eqref{eq:deficiency}: in that case
  $|\partial_N(S)| = d_N(v)$, while the values above and the assumptions on~$k$
  give $\delta(v) \le d_N(v)$. 
  If a connected set with at least two vertices contains at most one marked
  head~$v$ satisfying $d_N(v) \ge k-1$, its $\delta$-sum is at most~$1$: that
  head contributes at most~$2$ and, if present, has an ordinary neighbour
  contributing at most~$-1$, while all other vertices contribute at most~$0$.
  Thus~$S$ contains both marked heads, say~$b_1,b_2$, each satisfying
  $d_N(b_i) \ge k-1$, and $k \ge 4$.

  Let~$q$ be the number of ordinary vertices in~$S$.
  Since the two heads lie in the same bipartition class and $N[S]$ is
  connected, $q \ge 1$.
  Each head contributes at most~$2$ to the $\delta$-sum, and each ordinary
  vertex contributes at most~$-1$. 
  Hence~\eqref{eq:deficiency} gives $2 \le 4 - q - |\partial_N(S)|$.
  Thus either $q = 1$ and $|\partial_N(S)| \le 1$, or $q = 2$ and
  $|\partial_N(S)| = 0$.
  In the latter case, $S$ is the vertex set of a component of~$N$, both heads
  have even degree in~$N$, and both ordinary vertices have odd degree in~$N$.

  If $q = 1$, let~$a$ be the unique ordinary vertex in~$S$.
  Direct every edge of~$N[S]$ towards~$a$ and direct the possible boundary edge
  into~$S$.
  If the resulting indegree of~$a$ in~$N$ is odd, reverse one edge of~$N[S]$
  towards a head not incident with the boundary.
  Such an edge exists because~$N[S]$ is connected. 
  The indegree of~$a$ is then even and, after the prescribed arcs of~$P$ are
  restored, each marked head has total indegree~$1$ or~$2$.

  If $q = 2$, the underlying simple graph of~$N[S]$ has a perfect matching.
  Direct every edge of~$N[S]$ towards the ordinary vertices and reverse the
  direction of one edge corresponding to each edge of this matching. 
  Each ordinary vertex then has even indegree, since its degree in~$N$ is odd,
  and each marked head has one incoming edge from~$N[S]$ and one from~$P$,
  hence total indegree~$2$.

  There are no marked heads outside~$S$. 
  The possible boundary edge was directed into~$S$, so it contributes no
  indegree outside~$S$.
  Apply part~(ii) of \Cref{lem:parity-orientation} separately to the nontrivial
  components of $N[V(N) \setminus S]$, orienting their internal edges so that
  every indegree is even or~$1$.
  Together, these orientations and those already fixed inside~$S$, on the
  possible boundary edge, and on~$P$ give every vertex of~$M$ indegree even
  or~$1$ and give neither marked head indegree~$k$.
\end{proof}

\end{document}
